\section{Clean Release and Overcompression}

The contrast-concentration theorem establishes an exact finite
mechanism by which common-mode dominance can increase relative to
centered relational contrast.

It does not by itself determine an optimal loading regime.

This section introduces two internal concepts motivated by that
theorem:

\[
\text{clean release}
\]
and
\[
\text{overcompression}.
\]

Neither concept is part of the theorem proved in Sections 4--7.

\subsection{A declared output register}

Let
\[
\rho:X\longrightarrow H
\]
be a declared finite real-valued output register, with
\[
H\subseteq\mathbb R^n.
\]

For
\[
h=\rho(x),
\]
write the orthogonal decomposition
\[
h
=
P_0h+P_\perp h.
\]

The common component is
\[
P_0h
=
\overline h\,\one_n,
\]
while the centered component is
\[
P_\perp h
=
h-\overline h\,\one_n.
\]

Define
\[
C(h)
=
\|P_0h\|_2
\]
and
\[
R(h)
=
\|P_\perp h\|_2.
\]

A strong registered output need not maximize either quantity
separately.

A pure common-mode output may have large total magnitude while
carrying no internal coordinate contrast.

Conversely, a highly contrasted state need not be organized into a
robust outward registration.

\subsection{Coherence and individuation are different}


For the present internal mechanics, we distinguish two descriptive
properties.

\paragraph{Coherent registration.}
Multiple admissible contributions reinforce a common registered output
rather than cancelling one another.

\paragraph{Registered individuation.}
The declared output retains enough centered distinction to resolve the
relational alternatives relevant to that registration.

These properties need not vary together.

The upgraded native theorem now places an important restriction on any
proposed clean-release cycle.

Native relational loading itself does not suppress Paper 14 contrast
under the canonical completion register.

Instead,
\[
8\longrightarrow2\longrightarrow1
\]
sharpens the selected future and gives
\[
\zeta:
\frac{31}{55}
\longrightarrow
\frac{41}{65}
\longrightarrow
\frac{7}{10}.
\]

Thus any later loss of centered contrast requires a distinct
contrast-concentrating operation, such as the normalized quadratic,
or an explicitly proved composition law joining selection and
concentration.

Native selection and contrast suppression must not be silently
identified.

\subsection{Clean release}

\begin{definition}[Clean release]
A clean-release regime is a loading regime in which increasing
relational organization improves the coherence, robustness, or
magnitude of a declared outward registration while preserving the
independent contrast required by that output.
\end{definition}

The definition is projection-relative.

It does not assert a universal scalar measure of release quality.

A quantitative clean-release criterion would require at least:

\begin{enumerate}[label=(\roman*)]

\item a declared loading parameter;

\item a declared outward registration;

\item a measure of coherent registered output;

\item a measure of surviving independent contrast.

\end{enumerate}

No unique combination of these quantities is selected in this paper.

\subsection{Why maximum aggregate need not be maximum useful output}

Suppose a loading family is indexed by
\[
\lambda.
\]

Let
\[
C(\lambda)
\]
measure coherent outward registration and let
\[
D(\lambda)
\]
measure the independent contrast retained by that registration.

A useful output statistic
\[
U(\lambda)
\]
may depend positively on both.

For example, schematically,
\[
U(\lambda)
=
F(C(\lambda),D(\lambda)),
\]
where
\[
\frac{\partial F}{\partial C}>0,
\qquad
\frac{\partial F}{\partial D}>0.
\]

If
\[
C(\lambda)
\]
continues increasing while
\[
D(\lambda)
\]
begins decreasing, then maximum aggregate strength need not coincide
with maximum useful registered output.

An interior optimum may occur.

This is not guaranteed by the finite theorem.

It is a structural possibility.

\subsection{Overcompression}

\begin{definition}[Overcompression]
Overcompression begins when additional relational loading that
previously improved coherent registration instead reduces the
independent contrast required by the surviving output mode.
\end{definition}

The relevant transition is therefore not simply
\[
\text{large norm}
\longrightarrow
\text{larger norm}.
\]

It is
\[
\text{organization that preserves distinction}
\]
to
\[
\text{organization that suppresses required distinction}.
\]

Overcompression is thus defined by a change in the role of added
constraint or loading.

\subsection{Normalized quadratic model}

The normalized Thalean quadratic gives a simple limiting example.

For
\[
\widehat Q=\frac{Q}{98},
\]
the common component is fixed:
\[
P_0\widehat Q^{\,k}x
=
P_0x.
\]

The centered component contracts:
\[
R(\widehat Q^{\,k}x)
\le
\left(\frac{9}{49}\right)^k
R(x).
\]

Therefore indefinite normalized \(\widehat Q\)-iteration drives every
centered component asymptotically toward zero.

If a useful outward registration requires nonzero centered structure,
then sufficiently deep \(\widehat Q\)-iteration necessarily lies on
the overcompressed side of that requirement.

This conclusion is conditional on the declared output requiring such
contrast.

\subsection{Clean release as a boundary regime}


The upgraded finite theorem does not establish a single monotone
sequence from native loading to contrast saturation.

Instead, two separately proved operations exist:
\[
\text{native selection}
\]
and
\[
\text{quadratic contrast concentration}.
\]

A larger composed dynamics might contain a sequence of the schematic
form
\[
\text{relational organization}
\]
\[
\downarrow
\]
\[
\text{native selection}
\]
\[
\downarrow
\]
\[
\text{coherent registered output}
\]
\[
\downarrow
\]
\[
\text{clean-release regime}
\]
\[
\downarrow
\]
\[
\text{contrast-concentrating stage}
\]
\[
\downarrow
\]
\[
\text{overcompression}
\]
\[
\downarrow
\]
\[
\text{contrast saturation}.
\]

This remains an internal hypothesis.

Both native loading and quadratic contrast concentration are now
theorem-level finite operations.

What remains unproved is the dynamical law that composes them into one
cycle and determines whether an intermediate optimum exists.

\subsection{A possible diagnostic}

For a declared register with
\[
C(h)\neq0,
\]
define the relative contrast ratio
\[
\rho(h)
=
\frac{R(h)}{C(h)}.
\]

A clean-release regime may be sought where:

\[
C(h)
\]
is large or increasing,

while
\[
\rho(h)
\]
remains above the minimum required for the declared output.

Overcompression begins when additional loading causes
\[
\rho(h)
\]
to fall below that functional regime.

This is a diagnostic proposal.

No universal critical value of
\[
\rho
\]
is asserted.

\subsection{Relation to projection-relative saturation}


Clean release, native future saturation, and contrast saturation are
not synonyms.

Native registered-future saturation occurs when
\[
|\mathcal C(p)|=1.
\]

Contrast saturation concerns suppression of a declared centered
contrast statistic.

Clean release concerns a possible useful outward-registration regime
in a larger composed dynamics.

The present theorem does not establish a universal ordering among
these three notions.

In particular, native loading cannot itself be called overcompression
merely because it adds constraint: under the canonical completion
register it increases, rather than decreases, centered contrast.

Any transition from future sharpening to contrast suppression requires
an additional certified composition law.

\subsection{No physical promotion}

The terms in this section are internal relational-mechanics terms.

They are not identified here with:

\begin{itemize}

\item sonoluminescence emission;

\item hydrodynamic collapse;

\item gravitational collapse;

\item black-hole formation;

\item Hawking radiation;

\item plasma coherence;

\item or any physical emission mechanism.

\end{itemize}

Those would require separate empirical and physical crosswalks.

\subsection{Boundary}


The exact theorem establishes:

\begin{itemize}

\item native relational loading with
\[
8\to2\to1;
\]

\item retained contextual distinction at registered-future
saturation;

\item the exact retained-sector response
\[
\Delta^{\mathsf T}\Delta=\frac34H;
\]

\item and independent quadratic centered contraction under
\(\widehat Q\).

\end{itemize}

Clean release and overcompression remain conditional concepts because
no theorem yet specifies the dynamics that compose native selection
with quadratic concentration into one loading cycle.

Neither concept is promoted to physical interpretation here.
